\documentclass[10pt,a4paper]{amsart}
\usepackage{amsmath,amssymb,amsthm,graphicx}
% For commutative diagrams we use Xy-pic
\usepackage[all]{xy}

% We use 2cell for 2-commutative diagrams.
\xyoption{2cell}
% Theorem environments.
%
\theoremstyle{plain}
\newtheorem{theorem}[subsection]{Theorem}
\newtheorem{proposition}[subsection]{Proposition}
\newtheorem{lemma}[subsection]{Lemma}

\theoremstyle{definition}
\newtheorem{definition}[subsection]{Definition}
\newtheorem{example}[subsection]{Example}
\newtheorem{exercise}[subsection]{Exercise}
\newtheorem{situation}[subsection]{Situation}

\theoremstyle{remark}
\newtheorem{remark}[subsection]{Remark}
\newtheorem{remarks}[subsection]{Remarks}

\numberwithin{equation}{subsection}

% Macros
%
\def\lim{\mathop{\mathrm{lim}}\nolimits}
\def\colim{\mathop{\mathrm{colim}}\nolimits}
\def\Spec{\mathop{\mathrm{Spec}}}
\def\Hom{\mathop{\mathrm{Hom}}\nolimits}
\def\Ext{\mathop{\mathrm{Ext}}\nolimits}
\def\SheafHom{\mathop{\mathcal{H}\!\mathit{om}}\nolimits}
\def\SheafExt{\mathop{\mathcal{E}\!\mathit{xt}}\nolimits}
\def\Sch{\mathit{Sch}}
\def\Mor{\mathop{\mathrm{Mor}}\nolimits}
\def\Ob{\mathop{\mathrm{Ob}}\nolimits}
\def\Sh{\mathop{\mathit{Sh}}\nolimits}
\def\NL{\mathop{N\!L}\nolimits}
\def\CH{\mathop{\mathrm{CH}}\nolimits}
\def\proetale{{pro\text{-}\acute{e}tale}}
\def\etale{{\acute{e}tale}}
\def\QCoh{\mathit{QCoh}}
\def\Ker{\mathop{\mathrm{Ker}}}
\def\Im{\mathop{\mathrm{Im}}}
\def\Coker{\mathop{\mathrm{Coker}}}
\def\Coim{\mathop{\mathrm{Coim}}}

% Boxtimes
%
\DeclareMathSymbol{\boxtimes}{\mathbin}{AMSa}{"02}

%
% Macros for moduli stacks/spaces
%
\def\QCohstack{\mathcal{QC}\!\mathit{oh}}
\def\Cohstack{\mathcal{C}\!\mathit{oh}}
\def\Spacesstack{\mathcal{S}\!\mathit{paces}}
\def\Quotfunctor{\mathrm{Quot}}
\def\Hilbfunctor{\mathrm{Hilb}}
\def\Curvesstack{\mathcal{C}\!\mathit{urves}}
\def\Polarizedstack{\mathcal{P}\!\mathit{olarized}}
\def\Complexesstack{\mathcal{C}\!\mathit{omplexes}}
% \Pic is the operator that assigns to X its picard group, usage \Pic(X)
% \Picardstack_{X/B} denotes the Picard stack of X over B
% \Picardfunctor_{X/B} denotes the Picard functor of X over B
\def\Pic{\mathop{\mathrm{Pic}}\nolimits}
\def\Picardstack{\mathcal{P}\!\mathit{ic}}
\def\Picardfunctor{\mathrm{Pic}}
\def\Deformationcategory{\mathcal{D}\!\mathit{ef}}
\usepackage{cleveref}
\crefname{originalreference}{original reference}{original references}
\Crefname{originalreference}{Original reference}{Original references}
\makeatletter
\newlabel{proposition-topological-invariance}{{03SI}{source}{}{sourcepointer-1}{}}
\newlabel{proposition-topological-invariance@cref}{{[originalreference][0][]03SI}{[1][0][]source}{}{}{}}
\makeatother
\crefname{originalreference}{original reference}{original references}
\Crefname{originalreference}{Original reference}{Original references}
\makeatletter
\newlabel{theorem-vanishing-affine-curves}{{03SC}{source}{}{sourcepointer-2}{}}
\newlabel{theorem-vanishing-affine-curves@cref}{{[originalreference][0][]03SC}{[1][0][]source}{}{}{}}
\makeatother
\title{For a proper curve over a separably closed field cohomology is unchanged by extension of the base field}
\author{The Stacks Project Authors}
\date{stacks-project:etale-cohomology.tex:local-snapshot-20261011; GNU FDL 1.2 or later}
\begin{document}
\maketitle
\noindent\textit{Editorial note.} All delivered mathematical text is the author's own, copied from the pinned Stacks Project chapter file etale-cohomology.tex; no statement or proof step was rewritten, reordered or removed.\par\medskip
\noindent\textit{Editorial note.} The counted proof is the author's complete argument: the proof is written out in full: the square of fields, the pure inseparability of the extensions, the universality of the homeomorphisms, and the reduction to the algebraically closed case.\par\medskip
\noindent\textit{Editorial note.} The label lemma-base-change-dim-1-separably-closed names the counted statement, and the proof environment immediately following it in the source is delivered verbatim as its proof.\par\medskip
\noindent\textit{Editorial note.} The counted statement is carried by the author's own lemma environment in the chapter source.\par\medskip
\noindent\textit{Editorial note.} The statement and proof are copied from the pinned chapter file; the two cited inputs are kept as external pointers with their canonical Stacks tags.\par\medskip
\noindent\textit{Allowed standard prerequisites (editorial).}\begin{itemize}
\item The corresponding statement over algebraically closed fields. Same chapter, the vanishing and comparison theorem for affine curves over algebraically closed fields, which contains that case. Supplies the case of algebraically closed base fields.
\item Topological invariance of etale cohomology under universal homeomorphisms. Same chapter, the proposition identifying cohomology along purely inseparable extensions. Supplies the passage from separably closed to algebraically closed base fields.
\end{itemize}\medskip
\ifdefined\hypertarget\hypertarget{sourcepointer-1}{}\fi
\noindent\textit{Source reference.} Original reference 03SI: sources/stacks-etale-cohomology/pinned-original-source.tex, source lines 6548--6548. Prerequisite pointer invoked at proof line(s) 15798; it names a standard fact of the same chapter that the counted argument uses, and every step of the counted proof built on it is retained verbatim in the delivered body.\par\medskip
\ifdefined\hypertarget\hypertarget{sourcepointer-2}{}\fi
\noindent\textit{Source reference.} Original reference 03SC: sources/stacks-etale-cohomology/pinned-original-source.tex, source lines 15656--15656. Prerequisite pointer invoked at proof line(s) 15780; it names a standard fact of the same chapter that the counted argument uses, and every step of the counted proof built on it is retained verbatim in the delivered body.\par\medskip
\begin{lemma}
\label{lemma-base-change-dim-1-separably-closed}
Let $k'/k$ be an extension of separably closed fields.
Let $X$ be a proper scheme over $k$ of dimension $\leq 1$.
Let $\mathcal{F}$ be a torsion abelian sheaf on $X$.
Then the map $H^q_\etale(X, \mathcal{F}) \to
H^q_\etale(X_{k'}, \mathcal{F}|_{X_{k'}})$ is an isomorphism
for $q \geq 0$.
\end{lemma}

\begin{proof}
We have seen this for algebraically closed fields in
Theorem \ref{theorem-vanishing-affine-curves}.
Given $k \subset k'$ as in the statement of the lemma we can
choose a diagram
$$
\xymatrix{
k' \ar[r] & \overline{k}' \\
k \ar[u] \ar[r] & \overline{k} \ar[u]
}
$$
where $k \subset \overline{k}$ and $k' \subset \overline{k}'$ are
the algebraic closures. Since $k$ and $k'$ are separably closed
the field extensions
$\overline{k}/k$ and $\overline{k}'/k'$
are algebraic and purely inseparable. In this case the morphisms
$X_{\overline{k}} \to X$ and $X_{\overline{k}'} \to X_{k'}$
are universal homeomorphisms. Thus the cohomology of $\mathcal{F}$
may be computed on $X_{\overline{k}}$ and the cohomology
of $\mathcal{F}|_{X_{k'}}$ may be computed on $X_{\overline{k}'}$,
see Proposition \ref{proposition-topological-invariance}.
Hence we deduce the general case from the case of algebraically
closed fields.
\end{proof}

\end{document}
